% !TEX program = pdflatex
\documentclass[11pt]{article}
\usepackage[a4paper,margin=1in]{geometry}
\usepackage{amsmath,amssymb,amsfonts}
\usepackage{siunitx}
\usepackage{graphicx}
\usepackage{booktabs}
\usepackage{hyperref}
\usepackage{enumitem}
\usepackage{mathtools}
\usepackage{multirow}
\usepackage{float}
\usepackage{listings}
\usepackage{caption}
\usepackage{xcolor}

\hypersetup{colorlinks=true,linkcolor=blue,citecolor=blue,urlcolor=blue}
\lstset{basicstyle=\ttfamily\small,breaklines=true,columns=fullflexible,showstringspaces=false}

\title{Cardiovascular Risk Prediction: Clean Mathematical Specification and Validation Plan}
\author{Ryan O. van Gelder \\ \\ Citizen Gardens} \\ \\
\date{\today}

\begin{document}
\maketitle
\begin{abstract}
We specify, justify, and operationalize a reproducible pipeline for predicting major adverse cardiovascular events (MACE: myocardial infarction, stroke, cardiovascular death) at 1, 3, and 5 years, with both binary horizons and time-to-event formulations. We include cohort definitions, preprocessing, feature engineering from EHR and ECG, model families, calibration, evaluation, ablations, governance, and a validation and regulatory plan. Code snippets illustrate the training workflow.
\end{abstract}

\section{Objective}
Predict incident MACE in primary prevention using EHR and ECG. Outputs: calibrated risk at fixed horizons and individual survival curves. Decisions are supported by thresholds derived from explicit utility.

\section{Cohorts and Indexing}
\begin{itemize}[noitemsep]
  \item Adults (\(\geq 18\) y). Exclude prevalent CVD at baseline.
  \item Derive on Site A. Hold out Site B for external validation.
  \item Index date: first ECG with same-day labs and vitals within a \(\pm 90\)-day window.
\end{itemize}

\section{Notation}
For patient \(i\) and time index \(t\): observations \(y_{it}\in\mathbb{R}^{p}\), exogenous inputs \(u_{it}\in\mathbb{R}^{q}\), latent state \(x_{it}\in\mathbb{R}^{d}\). Event time \(\tilde T_i\), censoring \(C_i\), observed \(T_i=\min(\tilde T_i,C_i)\), indicator \(\Delta_i=\mathbf 1[\tilde T_i\le C_i]\).

\section{Data Preprocessing}
\begin{itemize}[noitemsep]
  \item Units harmonized to SI; winsorize at 0.5\% and 99.5\% after standardization.
  \item Missingness mask \(m_{it}\). Within-encounter forward fill; global-median imputation; record indicators. Sensitivity with MICE.
  \item Patient-level split: 70/10/20 train/val/test. External Site B held out.
  \item z-score normalization using training moments.
\end{itemize}

\section{Feature Engineering}
\subsection{Demographics, Vitals, Labs, Medications}
Age, sex, race/ethnicity; SBP, DBP, HR, BMI, SpO2; LDL-C, HDL-C, TC, TG, HbA1c, fasting glucose, creatinine/eGFR, hs-CRP; medication exposures (statins, antihypertensives, antidiabetics, antiplatelets) as binary indicators over 180 days.

\subsection{ECG Features and QA}
Intervals PR, QRS, QT and QTc (Bazett, Fridericia) \cite{Bazett1920,Fridericia1920}. Heart rate, axes, simple HRV metrics from 10 s strips when feasible. Extraction uses QRS detection (Pan--Tompkins) with wavelet-based delineation \cite{PanTompkins1985}. QA gates: per-lead SNR $>\,6$ dB; at least 8 good leads; morphology stability. Vendor harmonization maps XML/HL7 fields to a common schema. HRV from short strips is exploratory and excluded from primary models.

\subsection{Temporal Alignment and Recency}
Within a \(\pm 90\)-day window around the index ECG, choose nearest measurements. If multiple values exist, compute a recency-weighted snapshot with weight \(w=\exp(-\Delta t/30\,\text{days})\), and include \(\Delta t\) as a feature. If no value exists in-window, mark missing.

\section{Models}
\subsection{Static Baselines}
\paragraph{Logistic Regression at Horizon \(\tau\).} \(p_i(\tau)=\sigma(\beta_0^{(\tau)}+\beta^{(\tau)\top} z_i)\), with engineered index-time features \(z_i\). Trained with class weighting and elastic net.

\paragraph{Cox Proportional Hazards.} \(h(t\mid z_i)=h_0(t)\exp(\beta^\top z_i)\). Fit by maximizing partial likelihood \cite{Cox1972}. PH assumptions checked via Schoenfeld residuals \cite{Schoenfeld1982}. Time-varying extensions used if violated.

\subsection{Time-series Models}
\paragraph{Linear-Gaussian State Space.} Dynamics \(x_{i,t+1}=A x_{it}+B u_{it}+w_{it}\), observation \(y_{it}=C x_{it}+v_{it}\) with \(w_{it}\sim\mathcal N(0,Q)\) and \(v_{it}\sim\mathcal N(0,R)\). Parameters \(A,B,C,Q,R\) estimated by EM with Kalman smoothing \cite{ShumwayStoffer1982}. Filtered states feed a survival head.

\paragraph{Discrete-time Hazard (Dynamic).} For intervals \(k=1,\dots,K\) up to horizon \(\tau\):
\[\Pr(T_i=k\mid H_{i,k})=\sigma\big(\alpha_k+g(H_{i,k})\big),\]
where history \(H_{i,k}\) summarizes \((y_{i,1:k},u_{i,1:k})\) via GRU/Transformer or Kalman states.

\subsection{Competing Risks}
We model cause-specific hazards for MACE and non-CV death and report cumulative incidence via Aalen--Johansen \cite{AalenJohansen1978}. We include Fine--Gray subdistribution modeling as a secondary analysis \cite{FineGray1999}.

\section{Estimation and Calibration}
\subsection{Estimation}
Logistic models minimize cross-entropy with elastic net. Cox models maximize partial likelihood with L2 penalty. State-space learned by EM. Dynamic hazard trained with negative log-likelihood and early stopping.

\subsection{Calibration}
Binary horizons use isotonic regression \cite{ZadroznyElkan2002} or Platt scaling \cite{Platt1999}. Survival models assess calibration-in-the-large and slope, and use baseline-hazard recalibration when needed.

\section{Evaluation}
\begin{itemize}[noitemsep]
  \item Discrimination: AUROC and AUPRC at 1/3/5 years; Harrell's C for survival \cite{Harrell1982}.
  \item Accuracy: Brier and Integrated Brier \cite{Brier1950}.
  \item Calibration: reliability diagrams, expected calibration error, slope/intercept.
  \item Clinical utility: decision curve analysis \cite{VickersElkin2006}; net reclassification improvement vs baselines.
  \item Uncertainty: 1000 bootstrap replicates with patient clustering.
\end{itemize}

\section{Benchmarks}
ASCVD Pooled Cohort Equations \cite{Goff2014}, Framingham 2008 general CVD \cite{DAgostino2008}, QRISK3 \cite{HippisleyCox2017}, and SCORE2 where applicable \cite{SCORE22021}. Refit or recalibrate to the local cohort.

\section{Ablations and Stress Tests}
Comparisons: static logistic vs Cox vs dynamic hazard vs state-space; no-ECG vs with-ECG; labs-only vs labs+vitals vs +meds vs +ECG; imputation methods; regularization types; Kalman vs GRU vs Transformer summarizers; cross-site and temporal generalization; drift robustness.

\section{Decision Thresholds and Utility}
Choose thresholds by expected utility using cost matrix \(\{c_{FP},c_{FN},c_{TP},c_{TN}\}\). Report sensitivity, specificity, PPV, NPV at chosen operating points.

\section{Reproducibility and Governance}
Fixed seeds, versioned data, stored splits. Pre-registered metrics and endpoint. Model cards include intended use, calibration domain, and monitoring plan. Bias audits by sex, race/ethnicity, and age with equal-calibration checks.

\section{Monitoring and Model Updating}
Quarterly monitoring. Recalibration triggers: ECE $>0.03$ or calibration slope outside $[0.9,1.1]$. Annual refit with data through the prior quarter. Rolling predictions at each new encounter; alert cooldown 30 days.

\section{Validation and Regulatory Plan}
\subsection{ECG Feature Validation}
Reference set of $\geq 1000$ ECGs across devices and rhythms. Two cardiologists annotate fiducials and intervals; disagreements adjudicated. Report MAE, bias, ICC, and Bland--Altman limits. Acceptance: PR/QRS/QT MAE $\leq 8/6/12$ ms, QTc MAE $\leq 15$ ms.

\subsection{QA Failure Rates}
Track failure reasons: noise, lead loss, rhythm, vendor parse, short duration. KPI: overall QA fails $\leq 20\%$; per-device $\leq 25\%$.

\subsection{Medication Exposure and Adherence}
Daily exposure from dispensing with days-supply carryover and a 14-day grace period. Discontinuation sets exposure to zero. Half-life decay of 30 days while covered.

\subsection{Complexity ROI Gate}
Promote state-space only if it beats GRU hazard by AUROC $+0.02$ or C-index $+0.02$, Brier reduction $\geq 5\%$, and higher net benefit at clinical thresholds on internal and external tests.

\section{Minimal Dataset Schema}
\begin{table}[H]
  \centering
  \caption{Example schema files}
  \begin{tabular}{ll}
    \toprule
    File & Columns \\
    \midrule
    demographics.csv & id, index\_date, age, sex, race\_ethnicity \\
    vitals.csv & id, t, sbp, dbp, hr, bmi, spo2 \\
    labs.csv & id, t, ldl, hdl, tc, tg, hba1c, glucose, creatinine \\
    meds.csv & id, t, statin, ace\_arb, ccb, thiazide, beta\_blocker, metformin, sglt2, glp1 \\
    ecg.csv & id, t, hr, pr, qrs, qt, qtc, axis, rmssd, sdnn \\
    outcomes.csv & id, event\_time, event\_indicator, censor\_time \\
    \bottomrule
  \end{tabular}
\end{table}

\section{Pseudocode and Code Snippets}
\subsection*{Training Workflow}
\begin{lstlisting}[language=Python,caption={Training and evaluation skeleton}]
# Split
splits = split_patients(df, train=0.7, val=0.1, test=0.2)
scaler = fit_scaler(train)
X_train = featurize_index(train)
seq_train = featurize_sequences(train)

# Static models
log_reg = train_logistic(X_train, y_tau, class_weight='balanced', penalty='elasticnet')
cox = train_cox(X_train, (T, Delta), l2=1.0)

# Dynamic models (GRU hazard)
gru = train_gru_hazard(seq_train, (T, Delta), hidden=128, layers=1, dropout=0.1)

# Optional: state-space via EM + Kalman
a,b,c,q,r = em_kalman(seq_train)
z_train = kalman_filter_states(seq_train, a,b,c,q,r)
dhazard = train_discrete_hazard(z_train, (T, Delta))

# Calibration
cal_1y = fit_isotonic(val_pred_1y, val_lab_1y)

# Evaluation
evaluate_all(models=[log_reg, cox, gru], calibrators=[cal_1y], test=test)
external_validate(models=[log_reg, cox, gru], siteB=external)
\end{lstlisting}

\subsection*{Discrete-time Survival from Hazards}
\begin{lstlisting}[language=Python,caption={From interval hazards to survival and risk}]
# hazards: shape [n_patients, K]
S = np.cumprod(1.0 - hazards, axis=1)
# risk by tau: 1 - S[:, K_tau-1]
\end{lstlisting}

\subsection*{Recency-weighted Snapshot}
\begin{lstlisting}[language=Python,caption={Recency weighting within window}]
weights = np.exp(-delta_days / 30.0)
weighted_value = np.sum(values * weights) / np.sum(weights)
\end{lstlisting}

\section{Limitations}
This specification makes no causal claims. ECG HRV from 10 s strips is not used in primary models. All calibration and monitoring criteria must be validated in the target deployment environment.

\section{Fast Path to Proof}
Step 1: static logistic and Cox vs benchmarks. Step 2: calibration and decision curves. Step 3: GRU-based dynamic hazard. Step 4: external validation. Step 5: ship or kill based on net benefit and calibration.

\begin{thebibliography}{99}
\bibitem{Bazett1920} H.~C. Bazett. An analysis of the time-relations of electrocardiograms. \emph{Heart}, 7:353--370, 1920.
\bibitem{Fridericia1920} L.~S. Fridericia. Die Systolendauer im Elektrokardiogramm. \emph{Acta Medica Scandinavica}, 53:469--486, 1920.
\bibitem{PanTompkins1985} J.~Pan and W.~J. Tompkins. A real-time QRS detection algorithm. \emph{IEEE Trans. Biomed. Eng.}, 32(3):230--236, 1985.
\bibitem{Cox1972} D.~R. Cox. Regression models and life-tables. \emph{J. Royal Stat. Soc. B}, 34(2):187--220, 1972.
\bibitem{Schoenfeld1982} D. Schoenfeld. Partial residuals for the proportional hazards regression model. \emph{Biometrika}, 69(1):239--241, 1982.
\bibitem{ShumwayStoffer1982} R.~H. Shumway and D.~S. Stoffer. An approach to time series smoothing and forecasting using the EM algorithm. \emph{J. Time Series Analysis}, 3(4):253--264, 1982.
\bibitem{AalenJohansen1978} O. Aalen and S. Johansen. An empirical transition matrix for non-homogeneous Markov chains. \emph{Scand. J. Stat.}, 5(3):141--150, 1978.
\bibitem{FineGray1999} J.~P. Fine and R.~J. Gray. A proportional hazards model for the subdistribution of a competing risk. \emph{JASA}, 94(446):496--509, 1999.
\bibitem{ZadroznyElkan2002} B. Zadrozny and C. Elkan. Transforming classifier scores into accurate multiclass probability estimates. \emph{KDD}, 2002.
\bibitem{Platt1999} J. Platt. Probabilistic outputs for SVMs and comparisons to regularized likelihood methods. \emph{Advances in Large Margin Classifiers}, 1999.
\bibitem{Harrell1982} F.~E. Harrell Jr., R.~M. Califf, D.~B. Pryor, K.~L. Lee, and R.~A. Rosati. Evaluating the yield of medical tests. \emph{JAMA}, 247(18):2543--2546, 1982.
\bibitem{Brier1950} G.~W. Brier. Verification of forecasts expressed in terms of probability. \emph{Monthly Weather Review}, 78(1):1--3, 1950.
\bibitem{VickersElkin2006} A.~J. Vickers and E.~B. Elkin. Decision curve analysis: a novel method for evaluating prediction models. \emph{Med Decis Making}, 26(6):565--574, 2006.
\bibitem{Goff2014} D.~C. Goff et al. 2013 ACC/AHA guideline on the assessment of cardiovascular risk. \emph{Journal of the American College of Cardiology}, 63(25 Pt B):2935--2959, 2014.
\bibitem{DAgostino2008} R.~B. D'Agostino Sr. et al. General cardiovascular risk profile for use in primary care. \emph{Circulation}, 117(6):743--753, 2008.
\bibitem{HippisleyCox2017} J. Hippisley-Cox et al. Derivation and validation of QRISK3 risk prediction algorithms. \emph{BMJ}, 357:j2099, 2017.
\bibitem{SCORE22021} T. Piepoli et al. SCORE2 risk prediction algorithms: a new score for European populations. \emph{Eur Heart J}, 42(25):2439--2454, 2021.
\end{thebibliography}

\end{document}
